\section{Invitación a la comunidad}
% ── Frame: Tres frentes ──────────────────────────────────────────────────────
\begin{frame}{Súmate: tres frentes}
  \begin{columns}[T]
    \begin{column}{0.31\textwidth}
      % EDITAR
      \begin{exampleblock}{\footnotesize 1.\enspace Mapear lo que falta}
        \scriptsize El transporte periférico del Estado de México todavía tiene
        huecos en OSM: cada ruta que se mapea puede entrar a la red.
      \end{exampleblock}
    \end{column}
    \begin{column}{0.31\textwidth}
      % EDITAR
      \begin{block}{\footnotesize 2.\enspace Reportar y corregir}
        \scriptsize ¿Un trazo con errores o un sistema que falta? Se reporta en
        GitHub \textit{Discussions} para corregirlo o integrarlo.
      \end{block}
    \end{column}
    \begin{column}{0.31\textwidth}
      % EDITAR
      \begin{alertblock}{\footnotesize 3.\enspace Replicar en tu ciudad}
        \scriptsize Datos abiertos + OSM $\rightarrow$ API REST funciona donde haya
        GTFS y OSM: Bogotá, São Paulo, Buenos Aires…
      \end{alertblock}
    \end{column}
  \end{columns}

  \vspace{0.5cm}
  \centering
  % EDITAR: eco de la "Capa N" del bloque 5
  {\small\bfseries\color{SotmVerde} Y la Capa N: construye tu propio proyecto sobre la API.}
\end{frame}

% ── Frame: Ecosistema y enlaces ──────────────────────────────────────────────
\begin{frame}{Dónde encontrarnos}
  \begin{columns}[c]
    \begin{column}{0.56\textwidth}
      \centering
      \resizebox{\linewidth}{!}{% Diagrama: Ecosistema digital de Apimetro — cuatro plataformas de distribución
% Portal Web, Repositorio GitHub, Comunidad (Discussions) y Wiki
% Reemplaza las cuatro capturas de pantalla de la sección "El Futuro de Apimetro"
% Usado en: 4-Capas-Codigo/4-1-Apimetro/4-1-5-Futuro.tex  (fig:apimetro-ecosistema)
% Copia para la presentación SOTM (26 oct 2026) — no modifica la tesis:
%   · Wiki sin seed.sql (aún no se publica en GitHub)
\tikzsetnextfilename{apimetro-ecosistema}
\begin{tikzpicture}[
  central/.style={
    rectangle, draw=unamAzul, fill=unamAzul!18,
    text width=2.8cm, align=center,
    minimum height=1.5cm, rounded corners=5pt,
    font=\small\bfseries
  },
  plataforma/.style={
    rectangle, draw=unamAzul!55, fill=unamAzul!6,
    text width=3.2cm, align=center,
    minimum height=1.9cm, rounded corners=4pt,
    font=\footnotesize
  },
  flecha/.style={-Stealth, thick, color=unamAzul!55},
  etiq/.style={
    font=\scriptsize\itshape\color{gray!65},
    fill=white, inner sep=1.5pt, align=center
  }
]

%% — Nodo central —
\node[central] (api) at (0, 0) {
  \textbf{Apimetro}\\[3pt]
  {\scriptsize Motor de datos\\ZMVM}
};

%% — Satélites en disposición de cruz —
\node[plataforma] (portal) at (0, 3.6) {
  \textbf{Portal Web}\\
  \texttt{apimetro.dev}\\[3pt]
  {\scriptsize Documentación interactiva\\Swagger UI · estado del servidor}
};

\node[plataforma] (wiki) at (0, -3.6) {
  \textbf{Wiki}\\
  \texttt{GitHub Wiki}\\[3pt]
  {\scriptsize Guías de implementación\\Diccionario de datos}
};

\node[plataforma] (com) at (-5.0, 0) {
  \textbf{Comunidad}\\
  \texttt{Discussions}\\[3pt]
  {\scriptsize Propuestas · reportes\\Gobernanza abierta}
};

\node[plataforma] (repo) at (5.0, 0) {
  \textbf{Repositorio}\\
  \texttt{GitHub}\\[3pt]
  {\scriptsize Código fuente · Docker\\Licencia BSL~1.1}
};

%% — Flechas desde el nodo central —
\draw[flecha] (api.north) -- node[etiq, right=3pt] {interfaz pública} (portal.south);
\draw[flecha] (api.south) -- node[etiq, right=3pt] {implementación} (wiki.north);
\draw[flecha] (api.west)  -- node[etiq, above=3pt] {participación} (com.east);
\draw[flecha] (api.east)  -- node[etiq, above=3pt] {distribución} (repo.west);

\end{tikzpicture}
}\par
      \fuentefigura{Elaboración propia, 2026.}
    \end{column}
    \begin{column}{0.40\textwidth}
      \small
      % QR generados con el paquete qrcode desde la URL (sin imágenes externas);
      % fondo blanco con margen (zona de silencio) para que el lector los reconozca
      {\setlength{\fboxsep}{4pt}%
      \begin{block}{\footnotesize Apimetro}
        \begin{minipage}[c]{2.2cm}
          \colorbox{white}{\qrcode[height=1.75cm]{https://apimetro.dev}}
        \end{minipage}\hfill
        \begin{minipage}[c]{\dimexpr\linewidth-2.5cm}
          \raggedright\scriptsize Documentación, Swagger y API\\[2pt]
          \texttt{apimetro.dev}
        \end{minipage}
      \end{block}
      \begin{block}{\footnotesize VFTModel}
        \begin{minipage}[c]{2.2cm}
          \colorbox{white}{\qrcode[height=1.75cm]{https://github.com/galigaribaldi/VFTModel}}
        \end{minipage}\hfill
        \begin{minipage}[c]{\dimexpr\linewidth-2.5cm}
          \raggedright\scriptsize Repositorio (aún sin versión estable)\\[2pt]
          \texttt{github.com/}\\\texttt{galigaribaldi/VFTModel}
        \end{minipage}
      \end{block}}
    \end{column}
  \end{columns}
\end{frame}
